\lecture{6.6}{Variable Screening}
\label{lec:6-6}

\subsection{Why Screen Variables}

Variable screening is the process of selecting a subset of predictors before fitting a
learner. Screening can reduce variance, improve computation time, and allow the
library to include procedures representing different scientific or statistical choices.

In Super Learner, screening is itself part of the candidate algorithm. For example,
\texttt{SL.glm\_All} means a GLM using all variables, while
\texttt{SL.glm\_screen.corP} means a GLM fit only on variables selected by the
\texttt{screen.corP} screening rule.

\begin{keyidea}
Screening should be inside cross-validation. A screener must be fit on the training
fold only, then applied to determine which variables the learner uses for that fold.
Screening on the full dataset before cross-validation leaks validation information into
training.
\end{keyidea}

\subsection{Built-In Screeners}

Built-in screening functions can be listed with:

\begin{implementationbox}
\begin{lstlisting}[language=R]
listWrappers(what = "screen")
\end{lstlisting}
\end{implementationbox}

A screening function returns a logical vector indicating which columns of \texttt{X}
should be included. Its general structure is:

\begin{implementationbox}
\begin{lstlisting}[language=R]
screen.myScreen <- function(Y, X, family, obsWeights, id, ...) {
  # use training Y and X to decide which columns to keep
  keep <- rep(TRUE, ncol(X))
  return(keep)
}
\end{lstlisting}
\end{implementationbox}

\subsection{\texttt{screen.corP}}

The screener \texttt{screen.corP} selects variables based on marginal correlation tests
with the outcome. Conceptually, for each column $X_j$, it tests the association
between $X_j$ and $Y$ and keeps variables with small p-values. It also includes a
minimum-screening rule so that at least a small number of variables are retained.

This is a prediction-oriented screening rule. It is not a causal confounder-selection
procedure by itself. A variable can be a confounder even if its marginal correlation
with $Y$ is weak, and a variable can predict $Y$ without being needed for confounding
control.

\subsection{Combining Learners and Screeners}

To combine learners and screeners, pass \texttt{SL.library} as a list. Each element
contains a learner and a screener:

\begin{implementationbox}
\begin{lstlisting}[language=R]
sl_fit4 <- SuperLearner(
  Y = Y,
  X = data.frame(W1 = W1, W2 = W2, A = A),
  SL.library = list(c("SL.glm", "All"),
                    c("SL.glm", "screen.corP"),
                    c("SL.mean", "All")),
  family = gaussian(),
  method = "method.CC_LS",
  cvControl = list(V = 2),
  verbose = FALSE
)
\end{lstlisting}
\end{implementationbox}

The suffix in the printed learner name identifies the screener. A warning about
duplicate learners can occur when two learner-screener combinations lead to the same
selected variables and predictions.

\subsection{Custom Confounder Screening}

Sometimes an analyst wants a screener motivated by confounding control rather than
pure prediction. One simple historical rule is to include a covariate if adjusting for it
changes the treatment coefficient by more than a specified percentage. The rule below
always includes the treatment variable and screens other variables by the treatment
coefficient change criterion.

\begin{implementationbox}
\begin{lstlisting}[language=R]
screen_confounders <- function(Y, X, family, trt_name = "A",
                               change = 0.1, ...) {
  X <- as.data.frame(X)

  fit_init <- glm(as.formula(paste0("Y ~ ", trt_name)),
                  data = X, family = family)
  trt_coef <- coef(fit_init)[2]

  trt_col <- which(colnames(X) == trt_name)
  include <- rep(FALSE, ncol(X))
  include[trt_col] <- TRUE

  for (j in seq_len(ncol(X))[-trt_col]) {
    var_name <- colnames(X)[j]
    fit <- glm(as.formula(paste0("Y ~ ", trt_name, " + ", var_name)),
               data = X, family = family)
    new_trt_coef <- coef(fit)[2]
    include[j] <- abs((new_trt_coef - trt_coef) / trt_coef) > change
  }

  return(include)
}
\end{lstlisting}
\end{implementationbox}

\subsection{Treatment-Coefficient Change Criterion}

The change-in-estimate criterion compares
\[
\hat\beta_A^{\mathrm{unadjusted}}
\quad\text{and}\quad
\hat\beta_A^{\mathrm{adjusted\ for\ }X_j}.
\]
Variable $X_j$ is kept if
\[
\left|
\frac{\hat\beta_A^{\mathrm{adjusted\ for\ }X_j}
-\hat\beta_A^{\mathrm{unadjusted}}}
{\hat\beta_A^{\mathrm{unadjusted}}}
\right|>\delta,
\]
where $\delta$ might be $0.10$. This rule is heuristic. It depends on the outcome
model scale, can behave poorly when the unadjusted treatment coefficient is near zero,
and does not guarantee adjustment for all confounders.

\subsection{Interpreting Screened Learners}

A screened learner is a compound procedure: screening rule plus prediction algorithm.
For example, \texttt{SL.glm\_screen.corP} is not simply \texttt{SL.glm}; it is the
algorithm that first screens variables using \texttt{screen.corP} in each training fold
and then fits \texttt{SL.glm} using the selected variables.

\begin{assumptionbox}
Variable screening can remove variables needed for causal identification. If a
covariate is required for exchangeability based on subject-matter knowledge, it should
not be dropped solely because it has weak predictive association in the observed
sample. Prediction-oriented screening must be reconciled with the causal adjustment
set.
\end{assumptionbox}

\begin{estimationbox}
A practical compromise is to force known confounders into every nuisance model and
allow screening only among additional candidate predictors. This preserves the
scientific adjustment set while still allowing data-adaptive prediction improvements.
\end{estimationbox}
